% !TeX root = part7.tex
\documentclass[../final.tex]{subfiles}

\begin{document}

% Линии фейнмановских диаграмм: стрелка фермиона и волнистый фотон.
\tikzset{
    feyn fermion/.style={
        draw=rtext,
        line width=0.9pt,
        postaction={decorate},
        decoration={
            markings,
            mark=at position 0.55 with {\arrow{Latex}}
        }
    },
    feyn photon/.style={
        draw=rtext,
        line width=0.85pt,
        decorate,
        decoration={
            snake,
            amplitude=1.6pt,
            segment length=6pt
        }
    },
    feyn line/.style={
        draw=rtext,
        line width=0.9pt
    }
}

\practice{7}{15.10.2026}{Бесспиновая электродинамика. Кроссинг и переменные Мандельштама}

\rtheory{Пропагатор и вершина}

Напомним элементы диаграмм для бесспиновых частиц:
\begin{align*}
    &\text{пропагатор фотона ($s=1$, импульс $q$):}
    && \frac{-ig^{\mu\nu}}{q^2},\\
    &\text{вершинный множитель ($s=0$, $p\to p'$, фотон $s=1$):}
    && ie\,(p+p')^\mu.
\end{align*}

\task[1]{Бесспиновое \texorpdfstring{$e^-\mu^-$}{e-μ-}-рассеяние}
\label{task:spinless_emu_sem7}

\condition
Найти $d\sigma/d\Omega$ для рассеяния $e^-\mu^-\to e^-\mu^-$ бесспиновых частиц.

\solution
\subsection*{Амплитуда}

\begin{rbox}[left=16pt,right=16pt,top=12pt,bottom=10pt]{[]}
    \centering
    \begin{tikzpicture}[>=Latex,every node/.style={text=rtext}]
        \coordinate (V1) at (0,0.9);
        \coordinate (V2) at (0,-0.9);
        \draw[feyn fermion] (-1.6,2.3) -- (V1)
            node[pos=0.3,left] {$p_1$};
        \draw[feyn fermion] (V1) -- (1.6,2.3)
            node[pos=0.7,right] {$p_3$};
        \draw[feyn fermion] (-1.6,-2.3) -- (V2)
            node[pos=0.3,left] {$p_2$};
        \draw[feyn fermion] (V2) -- (1.6,-2.3)
            node[pos=0.7,right] {$p_4$};
        \draw[feyn photon] (V1) -- (V2);
        \draw[->,rtextmuted] (0.4,0.45) -- (0.4,-0.45)
            node[midway,right] {$q$};
        \node[above left] at (-1.6,2.3) {$e^-$};
        \node[below left] at (-1.6,-2.3) {$\mu^-$};
        \fill[rtext] (V1) circle(1.6pt);
        \fill[rtext] (V2) circle(1.6pt);
    \end{tikzpicture}
    \captionsetup{font=footnotesize,labelfont=bf,skip=5pt}
    \captionof{figure}{Диаграмма бесспинового рассеяния $e^-\mu^-\to e^-\mu^-$.}
    \label{fig:spinless_emu_sem7}
\end{rbox}

По правилам Фейнмана
\begin{equation*}
    -i\mathcal M=ie(p_1+p_3)_\mu
    \left(\frac{-ig^{\mu\nu}}{q^2}\right)
    ie(p_2+p_4)_\nu,
    \qquad
    q=p_1-p_3.
\end{equation*}
Отсюда
\begin{equation*}
    \mathcal M=-e^2\,\frac{1}{(p_1-p_3)^2}\,
    \bigl((p_1+p_3)\cdot(p_2+p_4)\bigr).
\end{equation*}

\subsection*{Кинематика в системе центра масс}

\rassumption{Работаем в системе центра масс и пренебрегаем массами}
(ультрарелятивистский случай). Тогда у всех частиц одна энергия $E$:
\begin{align*}
    p_1&=(E,\,E\vec n), & p_2&=(E,\,-E\vec n),\\
    p_3&=(E,\,E\vec n'), & p_4&=(E,\,-E\vec n'),
    \qquad \vec n\cdot\vec n'=\cos\theta.
\end{align*}
Знаменатель:
\begin{equation*}
    (p_1-p_3)^2=-2p_1\cdot p_3=-2E^2(1-\cos\theta).
\end{equation*}
Числитель:
\begin{align*}
    (p_1+p_3)\cdot(p_2+p_4)
    &=p_1p_2+p_1p_4+p_3p_2+p_3p_4\\
    &=2E^2+E^2(1+\cos\theta)+E^2(1+\cos\theta)+2E^2
    =2E^2(3+\cos\theta).
\end{align*}
Поэтому
\begin{equation*}
    \mathcal M=e^2\,\frac{3+\cos\theta}{1-\cos\theta}.
\end{equation*}

\subsection*{Сечение}

Для упругого рассеяния $p_f=p_i$, и по формуле
из задачи~1 занятия~6
\begin{equation*}
    \frac{d\sigma}{d\Omega}=\frac{1}{64\pi^2s}|\mathcal M|^2
    =\frac{e^4}{64\pi^2s}
    \left(\frac{3+\cos\theta}{1-\cos\theta}\right)^2,
    \qquad
    s=(E_1+E_2)^2.
\end{equation*}

\answer
С $e^2=4\pi\alpha$:
\begin{equation*}
    \frac{d\sigma}{d\Omega}=\frac{\alpha^2}{4s}
    \left(\frac{3+\cos\theta}{1-\cos\theta}\right)^2.
\end{equation*}

\task[2]{Рассеяние \texorpdfstring{$e^-\mu^+\to e^-\mu^+$}{e-μ+ → e-μ+}}
\label{task:spinless_emu_plus_sem7}

\solution
\begin{rbox}[left=16pt,right=16pt,top=12pt,bottom=10pt]{[]}
    \centering
    \begin{tikzpicture}[>=Latex,every node/.style={text=rtext}]
        \coordinate (V1) at (0,0.9);
        \coordinate (V2) at (0,-0.9);
        \draw[feyn fermion] (-1.6,2.3) -- (V1)
            node[pos=0.3,left] {$p_1$};
        \draw[feyn fermion] (V1) -- (1.6,2.3)
            node[pos=0.7,right] {$p_3$};
        \draw[feyn fermion] (V2) -- (-1.6,-2.3)
            node[pos=0.7,left] {$p_2$};
        \draw[feyn fermion] (1.6,-2.3) -- (V2)
            node[pos=0.3,right] {$p_4$};
        \draw[feyn photon] (V1) -- (V2);
        \draw[->,rtextmuted] (0.4,0.45) -- (0.4,-0.45)
            node[midway,right] {$q$};
        \node[above left] at (-1.6,2.3) {$e^-$};
        \node[below left] at (-1.6,-2.3) {$\mu^+$};
        \fill[rtext] (V1) circle(1.6pt);
        \fill[rtext] (V2) circle(1.6pt);
    \end{tikzpicture}
    \captionsetup{font=footnotesize,labelfont=bf,skip=5pt}
    \captionof{figure}{Рассеяние $e^-\mu^+\to e^-\mu^+$. На линии $\mu^+$
    стрелки направлены против импульсов $p_2$ и $p_4$.}
    \label{fig:spinless_emu_plus_sem7}
\end{rbox}

Вершина античастицы даёт множитель $ie(-p_2-p_4)$:
\begin{equation*}
    -i\mathcal M=ie(p_1+p_3)_\mu
    \left(\frac{-ig^{\mu\nu}}{q^2}\right)
    ie(-p_2-p_4)_\nu,
    \qquad
    q=p_1-p_3.
\end{equation*}
% В записи пропагатор здесь без минуса, i g^{μν}/q^2, а вершина μ+ записана как (+ie)(-p2-p4).
% Выше использован пропагатор из правил Фейнмана; общий знак на |M|^2 не влияет.
Амплитуда отличается от задачи~\ref{task:spinless_emu_sem7} только знаком,
поэтому \rresult{сечение то же самое}.

\task[3]{Аннигиляция \texorpdfstring{$e^+e^-\to\mu^+\mu^-$}{e+e- → μ+μ-}}
\label{task:spinless_annihilation_sem7}

\solution
Диаграмму с обменом фотоном между линиями $e$ и $\mu$, нарисованную
по аналогии с рассеянием, \rnote{рисовать так нельзя}: пара $e^+e^-$
должна исчезнуть, а пара $\mu^+\mu^-$ --- родиться. Правильная диаграмма ---
аннигиляция в виртуальный фотон.

\begin{rbox}[left=16pt,right=16pt,top=12pt,bottom=10pt]{[]}
    \centering
    \begin{tikzpicture}[>=Latex,every node/.style={text=rtext}]
        \coordinate (V1) at (-1.1,0);
        \coordinate (V2) at (1.1,0);
        \draw[feyn fermion] (V1) -- (-2.5,1.3)
            node[pos=0.8,above] {$e^+\ p_1$};
        \draw[feyn fermion] (-2.5,-1.3) -- (V1)
            node[pos=0.2,below] {$e^-\ p_2$};
        \draw[feyn fermion] (V2) -- (2.5,1.3)
            node[pos=0.8,above] {$\mu^-\ p_4$};
        \draw[feyn fermion] (2.5,-1.3) -- (V2)
            node[pos=0.2,below] {$\mu^+\ p_3$};
        \draw[feyn photon] (V1) -- (V2);
        \draw[->,rtextmuted] (-0.45,0.35) -- (0.45,0.35)
            node[midway,above] {$q$};
        \fill[rtext] (V1) circle(1.6pt);
        \fill[rtext] (V2) circle(1.6pt);
    \end{tikzpicture}
    \captionsetup{font=footnotesize,labelfont=bf,skip=5pt}
    \captionof{figure}{Аннигиляция $e^+e^-\to\mu^+\mu^-$ через виртуальный фотон.}
    \label{fig:spinless_annihilation_sem7}
\end{rbox}
% В записи μ+ нарисован вверху справа, μ- внизу; импульсы p3 (μ+) и p4 (μ-) те же.

Амплитуда:
\begin{equation*}
    -i\mathcal M=ie(-p_1+p_2)_\mu
    \left(\frac{-ig^{\mu\nu}}{q^2}\right)
    ie(-p_3+p_4)_\nu,
    \qquad
    q=p_1+p_2.
\end{equation*}

\rtheory{Кроссинг}

\rconcept{Кроссинг} --- метод получения амплитуды одного процесса из амплитуды
другого: античастица в начальном состоянии эквивалентна частице в конечном
с противоположным 4-импульсом, и наоборот.
\begin{equation*}
    \mathcal M(AB\to CD)\longrightarrow
    \mathcal M(\overline A,\overline B,\overline C,\overline D).
\end{equation*}

\begin{rbox}[left=16pt,right=16pt,top=12pt,bottom=10pt]{[]}
    \centering
    \begin{tikzpicture}[>=Latex,every node/.style={text=rtext}]
        \begin{scope}
            \draw[feyn line] (0,0) circle(0.75);
            \draw[rtextmuted,dashed] (0,1.15) -- (0,-1.15);
            \draw[feyn fermion] (-1.9,1.3) -- (-0.6,0.45)
                node[pos=0.2,above] {$e^-\ p_A$};
            \draw[feyn fermion] (-1.9,-1.3) -- (-0.6,-0.45)
                node[pos=0.2,below] {$\mu^-\ p_B$};
            \draw[feyn fermion] (0.6,0.45) -- (1.9,1.3)
                node[pos=0.8,above] {$e^-\ p_C$};
            \draw[feyn fermion] (0.6,-0.45) -- (1.9,-1.3)
                node[pos=0.8,below] {$\mu^-\ p_D$};
            \node at (0,-2.05) {$e^-\mu^-\to e^-\mu^-$};
        \end{scope}
        \begin{scope}[xshift=6.2cm]
            \draw[feyn line] (0,0) circle(0.75);
            \draw[rtextmuted,dashed] (0,1.15) -- (0,-1.15);
            \draw[feyn fermion] (-1.9,1.3) -- (-0.6,0.45)
                node[pos=0.2,above] {$e^-\ p_A'$};
            \draw[feyn fermion] (-0.6,-0.45) -- (-1.9,-1.3)
                node[pos=0.8,below] {$\mu^+\ p_B'$};
            \draw[feyn fermion] (0.6,0.45) -- (1.9,1.3)
                node[pos=0.8,above] {$e^-\ p_C'$};
            \draw[feyn fermion] (1.9,-1.3) -- (0.6,-0.45)
                node[pos=0.2,below] {$\mu^+\ p_D'$};
            \node at (0,-2.05) {$e^-\mu^+\to e^-\mu^+$};
        \end{scope}
        \begin{scope}[xshift=3.1cm,yshift=-4.9cm]
            \draw[feyn line] (0,0) circle(0.75);
            \draw[rtextmuted,dashed] (0,1.15) -- (0,-1.15);
            \draw[feyn fermion] (-1.9,1.3) -- (-0.6,0.45)
                node[pos=0.2,above] {$e^-\ p_A'$};
            \draw[feyn fermion] (-0.6,-0.45) -- (-1.9,-1.3)
                node[pos=0.8,below] {$e^+\ p_B'$};
            \draw[feyn fermion] (0.6,0.45) -- (1.9,1.3)
                node[pos=0.8,above] {$\mu^-\ p_C'$};
            \draw[feyn fermion] (1.9,-1.3) -- (0.6,-0.45)
                node[pos=0.2,below] {$\mu^+\ p_D'$};
            \node at (0,-2.05) {$e^-e^+\to\mu^-\mu^+$};
        \end{scope}
    \end{tikzpicture}
    \captionsetup{font=footnotesize,labelfont=bf,skip=5pt}
    \captionof{figure}{Кроссинг: из $e^-\mu^-\to e^-\mu^-$ получаем
    $e^-\mu^+\to e^-\mu^+$ и $e^-e^+\to\mu^-\mu^+$.
    Пунктир отделяет начальное состояние от конечного.}
    \label{fig:crossing_sem7}
\end{rbox}

Для перехода к рассеянию $e^-\mu^+\to e^-\mu^+$
\begin{equation*}
    p_B'=-p_D,\qquad p_C'=p_C,\qquad p_D'=-p_B.
\end{equation*}
Если же начальный $\mu^-$ перенести в конечное состояние (как $\mu^+$),
а конечный $e^-$ --- в начальное (как $e^+$), получаем аннигиляцию
$e^-e^+\to\mu^-\mu^+$ с
\begin{equation*}
    p_B'=-p_C,\qquad p_C'=+p_D,\qquad p_D'=-p_B.
\end{equation*}

\rtheory{Переменные Мандельштама}

Для процесса $AB\to CD$ вводятся \rconcept{переменные Мандельштама}:
\begin{align*}
    s&=(p_A+p_B)^2=(p_C+p_D)^2,\\
    t&=(p_A-p_C)^2=(p_D-p_B)^2,\\
    u&=(p_A-p_D)^2=(p_C-p_B)^2.
\end{align*}

\begin{rbox}[left=16pt,right=16pt,top=12pt,bottom=10pt]{[]}
    \centering
    \begin{tikzpicture}[>=Latex,every node/.style={text=rtext}]
        % s-канал
        \begin{scope}
            \coordinate (V1) at (-0.7,0);
            \coordinate (V2) at (0.7,0);
            \draw[feyn line] (-1.6,1.1) -- (V1) node[pos=0,above left] {$A$};
            \draw[feyn line] (-1.6,-1.1) -- (V1) node[pos=0,below left] {$B$};
            \draw[feyn line] (V2) -- (1.6,1.1) node[pos=1,above right] {$C$};
            \draw[feyn line] (V2) -- (1.6,-1.1) node[pos=1,below right] {$D$};
            \draw[feyn photon] (V1) -- (V2) node[midway,below=3pt] {$\mathcal M_s$};
            \node at (0,1.6) {$s$-канал};
        \end{scope}
        % t-канал
        \begin{scope}[xshift=4.6cm]
            \coordinate (V1) at (0,0.7);
            \coordinate (V2) at (0,-0.7);
            \draw[feyn line] (-1.1,1.6) -- (V1) node[pos=0,above left] {$\overline D$};
            \draw[feyn line] (V1) -- (1.1,1.6) node[pos=1,above right] {$C$};
            \draw[feyn line] (-1.1,-1.6) -- (V2) node[pos=0,below left] {$B$};
            \draw[feyn line] (V2) -- (1.1,-1.6) node[pos=1,below right] {$\overline A$};
            \draw[feyn photon] (V1) -- (V2) node[midway,right=3pt] {$\mathcal M_t$};
            \node at (0,2.25) {$t$-канал};
        \end{scope}
        % u-канал
        \begin{scope}[xshift=9.2cm]
            \coordinate (V1) at (0,0.7);
            \coordinate (V2) at (0,-0.7);
            \draw[feyn line] (-1.1,1.6) -- (V1) node[pos=0,above left] {$A$};
            \draw[feyn line] (V1) -- (1.1,-1.6) node[pos=1,below right] {$C$};
            \draw[feyn line] (-1.1,-1.6) -- (V2) node[pos=0,below left] {$\overline D$};
            \draw[feyn line] (V2) -- (1.1,1.6) node[pos=1,above right] {$\overline B$};
            \draw[feyn photon] (V1) -- (V2) node[midway,left=3pt] {$\mathcal M_u$};
            \node at (0,2.25) {$u$-канал};
        \end{scope}
    \end{tikzpicture}
    \captionsetup{font=footnotesize,labelfont=bf,skip=5pt}
    \captionof{figure}{Диаграммы $s$-, $t$- и $u$-каналов.}
    \label{fig:stu_channels_sem7}
\end{rbox}

Амплитуды каналов связаны кроссингом:
\begin{equation*}
    \mathcal M_t=\mathcal M_s(s\leftrightarrow t),
    \qquad
    \mathcal M_u=\mathcal M_s(s\leftrightarrow u)=\mathcal M_t(t\leftrightarrow u).
\end{equation*}

\subsection*{Сумма переменных}

\begin{equation*}
    s+t+u=m_A^2+m_B^2+m_C^2+m_D^2.
\end{equation*}
Действительно, раскрывая квадраты,
\begin{equation*}
    s+t+u=3m_A^2+m_B^2+m_C^2+m_D^2
    +2p_A\cdot(p_B-p_C-p_D),
\end{equation*}
а из сохранения импульса $p_B-p_C-p_D=-p_A$, так что последнее слагаемое
равно $-2m_A^2$.

\task[4]{Переменные Мандельштама для \texorpdfstring{$e^-e^-\to e^-e^-$}{e-e- → e-e-}}
\label{task:moller_mandelstam_sem7}

\condition
Для рассеяния $e^-e^-\to e^-e^-$ в системе центра масс заданы
масса $m$ и угол рассеяния $\theta$; $p=|\vec p_i|=|\vec p_f|$.
Найти $s$, $t$, $u$.

\solution
Электроны в конечном состоянии \rnote{мы не различаем}, поэтому рисуем
две диаграммы Фейнмана. \rnote{Если работаем с фермионами, то амплитуды
вычитаем}:
\begin{equation*}
    \mathcal M=\mathcal M_t-\mathcal M_u.
\end{equation*}

\begin{rbox}[left=16pt,right=16pt,top=12pt,bottom=10pt]{[]}
    \centering
    \begin{tikzpicture}[>=Latex,every node/.style={text=rtext}]
        \begin{scope}
            \coordinate (V1) at (0,0.55);
            \coordinate (V2) at (0,-0.55);
            \draw[feyn fermion] (-1.3,1.4) -- (V1) node[pos=0,above left] {$e^-\ A$};
            \draw[feyn fermion] (V1) -- (1.3,1.4) node[pos=1,above right] {$C$};
            \draw[feyn fermion] (-1.3,-1.4) -- (V2) node[pos=0,below left] {$e^-\ B$};
            \draw[feyn fermion] (V2) -- (1.3,-1.4) node[pos=1,below right] {$D$};
            \draw[feyn photon] (V1) -- (V2);
        \end{scope}
        \node at (2.6,0) {\Large$-$};
        \begin{scope}[xshift=5.2cm]
            \coordinate (V1) at (0,0.55);
            \coordinate (V2) at (0,-0.55);
            \draw[feyn fermion] (-1.3,1.4) -- (V1) node[pos=0,above left] {$A$};
            \draw[feyn fermion] (V1) -- (1.3,-1.4) node[pos=1,below right] {$D$};
            \draw[feyn fermion] (-1.3,-1.4) -- (V2) node[pos=0,below left] {$B$};
            \draw[feyn fermion] (V2) -- (1.3,1.4) node[pos=1,above right] {$C$};
            \draw[feyn photon] (V1) -- (V2);
        \end{scope}
        % Кинематика в с.ц.м.
        \begin{scope}[xshift=9.6cm]
            \draw[->,rtext,thick] (-1.3,0) -- (-0.15,0) node[pos=0,left] {$\vec p$};
            \draw[->,rtext,thick] (1.3,0) -- (0.15,0) node[pos=0,right] {$-\vec p$};
            \draw[->,p3,thick] (0,0) -- (1.2,0.9);
            \draw[->,p3,thick] (0,0) -- (-1.2,-0.9);
            \draw[rtextmuted] (0.6,0) arc[start angle=0,end angle=37,radius=0.6];
            \node at (0.85,0.25) {$\theta$};
        \end{scope}
    \end{tikzpicture}
    \captionsetup{font=footnotesize,labelfont=bf,skip=5pt}
    \captionof{figure}{Две диаграммы рассеяния $e^-e^-\to e^-e^-$ и кинематика
    в системе центра масс.}
    \label{fig:moller_sem7}
\end{rbox}

В системе центра масс
\begin{align*}
    p_A&=(E,\,\vec p), & p_B&=(E,\,-\vec p),\\
    p_C&=(E,\,\vec p'), & p_D&=(E,\,-\vec p'),
    \qquad |\vec p'|=p,\quad \vec p\cdot\vec p'=p^2\cos\theta,
\end{align*}
где $E^2=m^2+p^2$. Тогда
\begin{align*}
    s&=(2E)^2=4m^2+4p^2=4(m^2+p^2),\\
    t&=-(\vec p-\vec p')^2=-2p^2(1-\cos\theta),\\
    u&=-(\vec p+\vec p')^2=-2p^2(1+\cos\theta).
\end{align*}

\answer
\begin{equation*}
    s=4(m^2+p^2),\qquad
    t=-2p^2(1-\cos\theta),\qquad
    u=-2p^2(1+\cos\theta).
\end{equation*}
Проверка: $s+t+u=4m^2$, как и должно быть для четырёх электронов.

\end{document}
